\documentclass[11pt]{article}

% ---- Identidad documental (familia Rutsubo) ----
\usepackage[T1]{fontenc}
\usepackage[utf8]{inputenc}
\usepackage{mathpazo}                 % Palatino en el cuerpo
\usepackage[scaled=0.92]{helvet}      % Helvetica en encabezados
\usepackage[a4paper,margin=2.4cm,headsep=0.7cm]{geometry}
\usepackage{xcolor}
\usepackage{booktabs}
\usepackage{tabularx}
\usepackage{array}
\usepackage{enumitem}
\usepackage{titlesec}
\usepackage{fancyhdr}
\usepackage{microtype}
\usepackage[hidelinks]{hyperref}
\usepackage{caption}

% Paleta
\definecolor{grafito}{HTML}{26261F}
\definecolor{salvia}{HTML}{57634F}
\definecolor{salviacl}{HTML}{8A9880}
\definecolor{papel}{HTML}{F8F6F1}
\definecolor{linea}{HTML}{BDB8AC}
\definecolor{tinta}{HTML}{2B2B24}

\pagecolor{papel}
\color{tinta}
\captionsetup{labelfont={sf,bf,color=salvia},font=small}

% Encabezados sans, en salvia
\titleformat{\section}{\sffamily\large\bfseries\color{salvia}}{\thesection.}{0.5em}{}
\titleformat{\subsection}{\sffamily\normalsize\bfseries\color{grafito}}{\thesubsection}{0.5em}{}
\titlespacing{\section}{0pt}{1.4em}{0.6em}

% Regla fina de la familia
\newcommand{\regla}{\noindent\textcolor{linea}{\rule{\linewidth}{0.4pt}}}
\newcommand{\veredicto}[1]{\textbf{\textcolor{salvia}{#1}}}

% Pie/encabezado
\pagestyle{fancy}\fancyhf{}
\renewcommand{\headrulewidth}{0.4pt}
\renewcommand{\footrulewidth}{0pt}
\fancyhead[L]{\sffamily\footnotesize\textcolor{salvia}{Rutsubo}}
\fancyhead[R]{\sffamily\footnotesize\textcolor{linea}{Reporte de Métricas y Seguridad · Corte 1}}
\fancyfoot[C]{\sffamily\footnotesize\textcolor{linea}{\thepage}}
\renewcommand{\arraystretch}{1.25}

\emergencystretch=3em     % evita overfull sin patrones de guionado en español
\hyphenation{Light-house re-la-y dae-mon}

\begin{document}
\sloppy

% ============================ PORTADA ============================
\thispagestyle{empty}
\vspace*{2.2cm}
{\sffamily
\noindent\textcolor{salvia}{\rule{2.2cm}{2pt}}\\[1.4em]
\noindent{\Huge\bfseries\color{grafito} Primer Reporte de Métricas\\[0.25em] y Auditoría de Seguridad}\\[0.8em]
\noindent{\large\color{salvia} Rutsubo — agente de código local-first}\\[0.6em]
\regla\\[0.3em]
\noindent{\normalsize\color{linea} Corte 1 de 2 · Lighthouse + OWASP ZAP · evidencia reproducible}
}

\vfill
\noindent\renewcommand{\arraystretch}{1.35}
\begin{tabular}{@{}>{\sffamily\bfseries\color{salvia}}l l@{}}
Materia   & Tecnologías para el Desarrollo de Aplicaciones \\
Docente   & Mtro. Ali Santiago López Zunún \\
Autor     & Ramsés Alejandro Camas Nájera \\
Programa  & MCIT — Sistemas Inteligentes \\
Periodo   & 2.\textsuperscript{o} cuatrimestre \\
Corte     & 1 de 2 \\
Fecha     & 13 de julio de 2026 \\
\end{tabular}
\vspace*{1.2cm}
\newpage

% ============================ 1. RESUMEN ============================
\section{Resumen ejecutivo}
Este corte documenta la primera medición de rendimiento y auditoría de seguridad
de Rutsubo con evidencia reproducible. Se auditaron las \textbf{tres superficies
HTTP} del sistema conforme a los requisitos no funcionales: la superficie pública
estática (Astro) con \textbf{Lighthouse} (RNF-02, ADR-002) y las dos APIs —daemon
y relay— con \textbf{OWASP ZAP} (RNF-09).

\medskip
\noindent\textbf{Rendimiento (Lighthouse).} Las cuatro combinaciones
página\,$\times$\,perfil (Inicio y Login, en móvil y escritorio) obtienen
\veredicto{Performance 100, Accessibility 100, SEO 100 y Best Practices 96},
con First Contentful Paint entre 165 y 620~ms. El presupuesto RNF-02
(FCP $<$ 1\,500~ms) se \veredicto{cumple con holgura} en las cuatro. Un único
hallazgo (un 404 de \texttt{favicon.ico}) explica el 96 de Best Practices.

\medskip
\noindent\textbf{Seguridad (ZAP).} Cero hallazgos de riesgo Alto; cero fallos
(\textsf{FAIL}). El \veredicto{daemon no presentó ninguna alerta} sobre 11
endpoints autenticados: el endurecimiento de cabeceras de la fase API se sostuvo.
El relay concentra los \veredicto{2 hallazgos} del corte: una configuración CORS
permisiva (riesgo Medio) y la ausencia de la cabecera \texttt{X-Content-Type-Options}
(riesgo Bajo).

\medskip
\noindent\textbf{Correcciones.} Se comprometen \veredicto{3 correcciones} para el
corte 2 (una de Lighthouse, dos del relay); \veredicto{3 hallazgos} se descartan
con justificación técnica y \veredicto{1} se acepta temporalmente con su
disparador. La corrección en caliente prevista para fugas de información
(sección~4.3 de la orden) \veredicto{no fue necesaria}: los cuerpos de error de
ambas APIs son sobres \texttt{ErrorEnvelope} estructurados, sin trazas de pila.

\regla

% ============================ 2. METODOLOGÍA ============================
\section{Metodología}
\subsection{Herramientas y entorno}
\noindent
\begin{tabularx}{\linewidth}{@{}l X@{}}
\toprule
\textbf{Componente} & \textbf{Versión / detalle} \\
\midrule
Lighthouse (CLI)   & 13.4.0 (npm) \\
Navegador          & Google Chrome for Testing 149.0.7827.55 (headless) \\
OWASP ZAP          & 2.17.0 (imagen \texttt{zaproxy/zap-stable}) \\
Docker             & 29.1.3 \\
Sistema            & Ubuntu 24.04.3 LTS · WSL2 (kernel 6.6.87) \\
\bottomrule
\end{tabularx}

\subsection{Qué mide cada herramienta y por qué}
La medición de rendimiento se hace sobre la \textbf{superficie pública Astro}
—inicio y autenticación— por diseño: es la superficie renderizada de forma
estática (ADR-002, \emph{renderizado híbrido}) y la que fija el presupuesto de
carga de RNF-02. La aplicación (SPA con render en cliente) no produce
\emph{scores} de documento comparables y se reporta aparte solo como nota
metodológica. La auditoría de seguridad cubre las dos superficies HTTP expuestas
—daemon (\texttt{127.0.0.1:7431}) y relay (\texttt{127.0.0.1:8443})— tal como
exige RNF-09.

\subsection{Condiciones de medición (integridad del dato)}
Todos los números provienen de ejecuciones reales de esta sesión; los artefactos
crudos (JSON) se conservan en \texttt{docs/audits/corte1/} con su hash SHA-256
(anexo). Condiciones:
\begin{itemize}[leftmargin=1.2em,itemsep=1pt]
  \item Lighthouse se corre contra el \textbf{build de producción} (\texttt{astro build})
        servido de forma estática, nunca el servidor de desarrollo.
  \item \textbf{3 corridas} por combinación; se reporta la \textbf{mediana} del
        \emph{score} de Performance (Lighthouse tiene varianza). Las tres corridas
        de cada combinación coincidieron.
  \item ZAP se ejecuta como \textbf{proxy de escaneo pasivo} sobre los endpoints
        enumerados de los contratos C-1/C-2. En una API JSON sin raíz navegable el
        \emph{spider} clásico no descubre rutas; conducir peticiones reales a través
        del proxy audita las respuestas efectivas (200 autenticadas, 401, 404). El
        daemon se auditó \textbf{con el \emph{Bearer} local}; sus endpoints
        protegidos respondieron 200. Las instancias auditadas son \textbf{efímeras y
        desechables} (data-dir temporal y \emph{workspace} de prueba con archivos
        dummy), nunca trabajo real.
\end{itemize}

\subsection{Comandos ejecutados}
Reproducibles sobre el mismo build y las mismas instancias efímeras.

\medskip
\noindent\textit{Superficie pública (build de producción servido estático):}
\begin{quote}\ttfamily\scriptsize\raggedright
npm run build\quad(astro)\ \ →\ \ python3 -m http.server 4321 (dist/)
\end{quote}

\noindent\textit{Lighthouse — 4 combinaciones $\times$ 3 corridas (mediana):}
\begin{quote}\ttfamily\scriptsize\raggedright
CHROME\_PATH=<chromium> lighthouse \\
\ \ http://127.0.0.1:4321/\{,login/\} [-{}-preset=desktop] \\
\ \ -{}-output=json,html \ -{}-chrome-flags="-{}-headless=new ..."
\end{quote}

\noindent\textit{OWASP ZAP — escaneo pasivo por proxy (contenedor \texttt{zaproxy/zap-stable}):}
\begin{quote}\ttfamily\scriptsize\raggedright
docker run -d -p 8090:8090 zaproxy/zap-stable zap.sh -daemon -port 8090 \\
\ \ -config api.disablekey=true \\
curl -x http://127.0.0.1:8090 -H "Authorization: Bearer <token>" \\
\ \ http://host.docker.internal:7431/v1/\{health,sessions,approvals,...\} \\
curl "http://127.0.0.1:8090/OTHER/core/other/jsonreport/" -o zap-daemon.json
\end{quote}

\noindent Nota de red: en Docker Desktop sobre WSL2 el contenedor alcanza el host
por \texttt{host.docker.internal} (no por \texttt{-{}-network=host}); es el mismo
daemon/relay de \texttt{127.0.0.1}, comprobado por respuesta idéntica.

\subsection{WebPageTest}
La rúbrica lo admite como alternativa a Lighthouse. No se usa: exige una URL
pública y este corte se mide en \texttt{localhost} sin despliegue. Lighthouse
cubre íntegramente el requisito RNF-02.

\regla

% ============================ 3. LIGHTHOUSE ============================
\section{Resultados de Lighthouse}
Valores de la corrida mediana por combinación. \emph{Scores} en 0--100; métricas
de laboratorio en milisegundos (CLS es adimensional).

\medskip
\begin{center}\footnotesize
\begin{tabular}{@{}l l *{9}{c}@{}}
\toprule
\textbf{Página} & \textbf{Perfil} & \textbf{Perf} & \textbf{A11y} & \textbf{BP} & \textbf{SEO} & \textbf{FCP} & \textbf{LCP} & \textbf{TBT} & \textbf{CLS} & \textbf{SI} \\
\midrule
Inicio & mobile & 100 & 100 & 96 & 100 & 618 & 751 & 0 & 0.000 & 618 \\
Inicio & desktop & 100 & 100 & 96 & 100 & 165 & 201 & 0 & 0.000 & 165 \\
Login & mobile & 100 & 100 & 96 & 100 & 620 & 751 & 0 & 0.000 & 620 \\
Login & desktop & 100 & 100 & 96 & 100 & 166 & 201 & 0 & 0.000 & 166 \\
\bottomrule
\end{tabular}
\end{center}
\normalsize

\subsection{Veredicto contra el presupuesto RNF-02}
\noindent
\begin{tabularx}{\linewidth}{@{}l l l@{}}
\toprule
\textbf{Umbral (RNF-02 / metas del proyecto)} & \textbf{Observado} & \textbf{Resultado} \\
\midrule
FCP $<$ 1\,500~ms & 165--620~ms (máx.\ 620) & \veredicto{Cumple} \\
Performance $\ge$ 90 & 100 en las 4 & \veredicto{Cumple} \\
Accessibility $\ge$ 95 & 100 en las 4 & \veredicto{Cumple} \\
\bottomrule
\end{tabularx}

\subsection{Oportunidades y diagnósticos}
Lighthouse \textbf{no reportó oportunidades de rendimiento con ahorro estimado}:
la superficie estática ya es óptima (TBT 0~ms, CLS 0, sin recursos que diferir).
El único punto por debajo de 100 es \textbf{Best Practices = 96}, causado por
\texttt{errors-in-console}: el navegador registra un error de red \texttt{404} al
solicitar \texttt{/favicon.ico}, recurso que el sitio no declara. Es el hallazgo
H1 de la tabla de correcciones.

\subsection{Nota sobre la aplicación (SPA)}
La aplicación autenticada es de \emph{render en cliente} (CSR, ADR-002): sirve un
documento mínimo y construye la interfaz en el navegador. Un \emph{score} de
documento de Lighthouse sobre ella no es comparable con el de una página estática
—penaliza el arranque del \emph{bundle} de forma que no refleja la experiencia
real tras la hidratación— por lo que \textbf{no se reporta como métrica de
documento} en este corte, en línea con RNF-02, que fija el presupuesto sobre la
superficie pública. La SPA se gobierna por métricas de aplicación (tiempo a
interactivo tras autenticación, latencia de \emph{stream} de eventos C-3) que se
instrumentan en la fase de observabilidad, no con Lighthouse.

\regla

% ============================ 4. ZAP ============================
\section{Resultados de OWASP ZAP}
Resumen por superficie. La cuenta de alertas es por \emph{tipo}; entre paréntesis,
el número de instancias.

\medskip
\noindent
\begin{tabularx}{\linewidth}{@{}l c c X@{}}
\toprule
\textbf{Superficie} & \textbf{URLs auditadas} & \textbf{Alertas} & \textbf{Riesgo} \\
\midrule
Daemon (\texttt{7431}, autenticado) & 11 & 0 & --- \\
Relay (\texttt{8443}) & 6 & 2 & 1 Medio, 1 Bajo \\
\bottomrule
\end{tabularx}

\subsection{Daemon — sin hallazgos}
Sobre 11 endpoints autenticados (\texttt{/v1/health}, \texttt{/v1/sessions},
\texttt{/v1/approvals}, \texttt{/v1/rules}, \texttt{/v1/config/*},
\texttt{/v1/audit}) ZAP no levantó ninguna alerta. La captura directa de cabeceras
sobre las respuestas 200 confirma el endurecimiento de la fase API; el relay, en
contraste, no aplica todavía esta capa (origen de H3):

\medskip
\noindent
\begin{tabular}{@{}l l l@{}}
\toprule
\textbf{Cabecera de seguridad} & \textbf{Daemon} & \textbf{Relay} \\
\midrule
\texttt{X-Content-Type-Options} & \texttt{nosniff} & (ausente) \\
\texttt{X-Frame-Options} & \texttt{DENY} & (ausente) \\
\texttt{Referrer-Policy} & \texttt{no-referrer} & (ausente) \\
\texttt{Cache-Control} & \texttt{no-store} & (ausente) \\
\texttt{Server} (versión) & (ausente, sin fuga) & (ausente, sin fuga) \\
\texttt{Access-Control-Allow-Origin} & (no expuesto) & \texttt{*} (origen de H2) \\
\bottomrule
\end{tabular}

\subsection{Relay — 2 hallazgos}
\noindent
\begin{tabular}{@{}p{5.6cm} l c l@{}}
\toprule
\textbf{Alerta (ZAP)} & \textbf{Riesgo} & \textbf{Inst.} & \textbf{Clasificación} \\
\midrule
Cross-Domain Misconfiguration \newline {\footnotesize(id 10098, CWE-264, WASC-14)} & Medio & 5 & Corregir (H2) \\[0.3em]
X-Content-Type-Options Header Missing \newline {\footnotesize(id 10021, CWE-693, WASC-15)} & Bajo & 1 & Corregir (H3) \\
\bottomrule
\end{tabular}

\medskip
\noindent El hallazgo Medio corresponde a la cabecera \texttt{Access-Control-Allow-Origin: *}
que el relay emite en todas sus respuestas; el Bajo, a la ausencia de
\texttt{nosniff} en el relay (cabecera que el daemon sí incluye).

\regla
\newpage

% ============================ 5. CORRECCIONES ============================
\section{Análisis de hallazgos y correcciones}
Cada hallazgo se clasifica en \textbf{corregir} (real, con fix concreto),
\textbf{no corregir con justificación} (falso positivo o no aplicable por
arquitectura) o \textbf{aceptado temporalmente} (real pero pospuesto, con el
disparador que lo reactiva).

\medskip
\noindent\footnotesize
\begin{tabularx}{\linewidth}{@{}p{0.6cm} X l l@{}}
\toprule
\textbf{\#} & \textbf{Hallazgo · análisis y decisión} & \textbf{Riesgo} & \textbf{Req.} \\
\midrule

H1 & \textbf{404 de \texttt{/favicon.ico} (Lighthouse, BP 96).} El sitio Astro no
declara favicon; el navegador lo pide de forma implícita y falla.
\newline\textit{$\rightarrow$ Corregir (corte 2):} añadir \texttt{public/favicon.svg}
y \texttt{<link rel="icon">} en el layout base. Lleva Best Practices a 100. &
Bajo & RNF-02 \\
\addlinespace

H2 & \textbf{CORS permisivo en el relay (\texttt{Allow-Origin: *}, ZAP Medio).}
Proviene de \texttt{CorsLayer::permissive()}, necesario en M2 para que la SPA
(origen distinto) alcance el REST del relay entre orígenes.
\newline\textit{$\rightarrow$ Corregir (corte 2):} restringir a una lista blanca
configurable (\texttt{RELAY\_SPA\_ORIGINS}); permisivo solo bajo bandera de
desarrollo. &
Medio & RNF-09 \\
\addlinespace

H3 & \textbf{Falta \texttt{X-Content-Type-Options} en el relay (ZAP Bajo).} El
daemon ya fija \texttt{nosniff}; el relay no aplica la misma capa de cabeceras.
\newline\textit{$\rightarrow$ Corregir (corte 2):} añadir al relay la capa de
seguridad del daemon (\texttt{nosniff}, \texttt{X-Frame-Options},
\texttt{Referrer-Policy}, \texttt{Cache-Control: no-store}). &
Bajo & RNF-09 \\
\addlinespace

H4 & \textbf{Sin Content-Security-Policy en las respuestas de las APIs.} La CSP
protege documentos HTML frente a inyección de scripts; ambas superficies sirven
\texttt{application/json}. La CSP de la SPA vive en su propio \texttt{index.html}.
\newline\textit{$\rightarrow$ No corregir (justificado):} no aplica a una API JSON. &
Info & --- \\
\addlinespace

H5 & \textbf{Sin HSTS / TLS en las APIs locales.} El daemon acepta solo conexiones
de \emph{loopback} (RNF-04) y no usa TLS localmente; el relay corre HTTP en LAN
sin despliegue público (ADR-009: el buzón en claro mantiene el relay dentro de
LAN hasta M3).
\newline\textit{$\rightarrow$ Aceptado temporalmente.} \textit{Disparador:} el
despliegue público con certificado del \emph{hosting} lo convierte en corrección
obligatoria. &
Bajo & RNF-04 \\
\addlinespace

H6 & \textbf{\texttt{/v1/health} sin autenticación.} Es una sonda de \emph{liveness};
devuelve solo estado y versión, ningún dato sensible.
\newline\textit{$\rightarrow$ No corregir (justificado):} acceso anónimo deliberado
por diseño. &
Info & --- \\
\addlinespace

H7 & \textbf{El 400 del daemon repite el mensaje del \emph{parser} JSON.} Ante un
cuerpo malformado devuelve el error de sintaxis (p.\,ej.\ ``\emph{key must be a
string at line 1 column 2}''), sin traza de pila ni estado interno.
\newline\textit{$\rightarrow$ No corregir (justificado):} ayuda al cliente y no
filtra información sensible; el sobre de error es \texttt{ErrorEnvelope}
estructurado. &
Info & --- \\
\bottomrule
\end{tabularx}
\normalsize

\medskip
\noindent\textbf{Corrección en caliente (sección 4.3).} No se aplicó ninguna: la
condición que la habilita —fugas de trazas de pila o cuerpos crudos del
\emph{upstream} en los errores— no se observó. Ambas APIs devuelven
\texttt{\{"error":\{"code","message"\}\}} sin filtración, en línea con el criterio
de terminación de la fase API.

\regla

% ============================ 6. PLAN CORTE 2 ============================
\section{Plan hacia el corte 2}
Correcciones comprometidas, cada una con su criterio de verificación objetivo
(re-medición o re-escaneo que demuestre el delta):

\begin{enumerate}[leftmargin=1.4em,itemsep=3pt]
  \item \textbf{Favicon del sitio público (H1).} Verificación: re-correr Lighthouse
        en las 4 combinaciones; \texttt{errors-in-console} sin instancias y
        \textbf{Best Practices = 100}.
  \item \textbf{Lista blanca de orígenes CORS en el relay (H2).} Verificación:
        re-escaneo ZAP del relay; la alerta \emph{Cross-Domain Misconfiguration}
        (10098) desaparece con un origen no autorizado y persiste el acceso del
        origen configurado de la SPA.
  \item \textbf{Capa de cabeceras de seguridad en el relay (H3).} Verificación:
        re-escaneo ZAP; la alerta \emph{X-Content-Type-Options} (10021) desaparece;
        el \emph{delta} de cabeceras daemon\,$\leftrightarrow$\,relay queda en cero.
\end{enumerate}

\noindent Los hallazgos H4, H6 y H7 no requieren acción; H5 permanece aceptado
hasta el despliegue público (M3), momento en que TLS/HSTS pasan a ser obligatorios.

\regla

% ============================ 7. ANEXO ============================
\section{Anexo — evidencia cruda e integridad}
Los artefactos crudos que respaldan cada cifra viven en el repositorio del núcleo,
bajo \texttt{docs/audits/corte1/}: los JSON+HTML de Lighthouse
(\texttt{lighthouse/}) y de ZAP (\texttt{zap/}). A continuación, el hash SHA-256
de cada JSON (prefijo de 32 hex) para verificar que las cifras del reporte no se
alteraron.

\medskip
\noindent\footnotesize
\begin{tabular}{@{}p{7.2cm} l@{}}
\toprule
\textbf{Artefacto (JSON crudo)} & \textbf{SHA-256 (prefijo)} \\
\midrule
\texttt{\scriptsize lighthouse/landing-desktop-run1.report.json} & \texttt{\scriptsize cefc3d811a9fe2d09cef1e8470eddde6}\ldots \\
\texttt{\scriptsize lighthouse/landing-desktop-run2.report.json} & \texttt{\scriptsize e256221aee4e6ca3821d828b5aa4eaa4}\ldots \\
\texttt{\scriptsize lighthouse/landing-desktop-run3.report.json} & \texttt{\scriptsize b58f7a8925d3e34aa0fbc9618e2346fe}\ldots \\
\texttt{\scriptsize lighthouse/landing-mobile-run1.report.json} & \texttt{\scriptsize 50f921c4c7001c77748faaf38e42521f}\ldots \\
\texttt{\scriptsize lighthouse/landing-mobile-run2.report.json} & \texttt{\scriptsize 3decbc4c18a30f6c6f7fdd523753aeb2}\ldots \\
\texttt{\scriptsize lighthouse/landing-mobile-run3.report.json} & \texttt{\scriptsize 3ae254a34f79af4b954fd65baa4f71c6}\ldots \\
\texttt{\scriptsize lighthouse/login-desktop-run1.report.json} & \texttt{\scriptsize b9cf94a76f9a1fdcbcdec10426059275}\ldots \\
\texttt{\scriptsize lighthouse/login-desktop-run2.report.json} & \texttt{\scriptsize 4cc5ed56341339506dced32363516a1b}\ldots \\
\texttt{\scriptsize lighthouse/login-desktop-run3.report.json} & \texttt{\scriptsize c91d680f7086535ca6c1c00e5bdb129c}\ldots \\
\texttt{\scriptsize lighthouse/login-mobile-run1.report.json} & \texttt{\scriptsize bdaf32cae991cb6342b535d40f719783}\ldots \\
\texttt{\scriptsize lighthouse/login-mobile-run2.report.json} & \texttt{\scriptsize 7f0b3e0c714e95c2c1b2953e28dbfa63}\ldots \\
\texttt{\scriptsize lighthouse/login-mobile-run3.report.json} & \texttt{\scriptsize 4f88a51546d335d04388cc654f572e30}\ldots \\
\texttt{\scriptsize zap/zap-daemon.json} & \texttt{\scriptsize 5c1d06a319b357325bf270e6d59ad2a7}\ldots \\
\texttt{\scriptsize zap/zap-relay.json} & \texttt{\scriptsize 47178d2310408659e2b526be53b699b2}\ldots \\
\bottomrule
\end{tabular}
\normalsize

\medskip
\regla
\vspace{0.4em}
\noindent{\footnotesize\textcolor{linea}{Rutsubo · Reporte de Métricas y
Auditoría de Seguridad · Corte 1 de 2 · Ramsés Alejandro Camas Nájera · MCIT-UPCh
· julio 2026. Toda métrica proviene de una ejecución real de esta sesión; los
artefactos crudos y sus hashes acompañan este documento.}}

\end{document}
